\chapter{Literature Review}
\label{chap:litreview}

This chapter reviews the work the project builds on. It covers personalised
federated learning, clustered and multi-tier variants, applications to
intrusion detection, and the data quality literature specific to CICIDS2017.
The final section identifies the gap the project addresses.

\section{Federated learning under heterogeneity}
\label{sec:fl-hetero}

Federated averaging \parencite{mcmahan-2017} trains a single global model by
averaging client updates. It assumes clients hold data drawn from a similar
distribution. When they do not, the averaged model degrades, and the degradation
is worse the more the client distributions differ.

Two families of response exist. Personalised methods give each client its own
model while sharing structure with the others. Clustered methods group clients
by similarity and share within groups rather than globally.

pFedMe \parencite{dinh-2020} takes the first route using Moreau envelopes. Each
client solves a regularised local problem whose solution is pulled toward the
global model by a penalty term, so the client model can specialise without
drifting arbitrarily far. Ditto \parencite{li-2021} and Per-FedAvg
\parencite{fallah-2020} pursue related ideas through different regularisation
and meta-learning respectively.

\section{Clustered and multi-tier methods}
\label{sec:clustered}

Clustered federated learning assumes clients fall into groups. Sattler et al.
\parencite{sattler-2020} recursively bipartition clients by the cosine
similarity of their updates, splitting a cluster when its members' gradients
point in sufficiently different directions. The method splits but never merges
or reassigns, which limits it when the grouping is initially wrong.

FlexCFL \parencite{duan-2021} adds client migration between clusters to handle
distribution shift, keeping the number of clusters fixed. FedDrift
\parencite{jothimurugesan-2023} handles concept drift with both splits and
merges. A recent survey \parencite{cfl-survey-2025} organises these approaches
and notes that basic clustered methods, limited to splitting, perform poorly
under complex drift.

HCFL+ \parencite{hcfl-2025} generalises clustered federated learning into a
four-tier framework with dynamic re-clustering, unifying soft and hard
clustering under a common objective.

SCMoE-PFL \parencite{scmoe-2026} combines soft clustering with a
mixture-of-experts. Its multi-centre threshold clustering allows a client to
belong to several clusters at once. Each cluster produces an expert model, each
client separately trains a private model, and an energy-aware gating network
fuses them. The clustering signal is the client gradient, L2-normalised to
discard magnitude and projected by PCA before cosine similarity is taken, the
projection being necessary because cosine similarity becomes uninformative in
high dimensions.

PerMFL \parencite{permfl-2024} differs from all of these in providing
convergence guarantees. It also differs in assuming the grouping is given
rather than derived. Its device update is
\begin{equation}
\theta_{i,j}^{t,k,l+1} = \theta_{i,j}^{t,k,l} - \alpha\nabla f_{i,j}(\theta_{i,j}^{t,k,l}) - \alpha\lambda\left(\theta_{i,j}^{t,k,l} - w_i^{t,k}\right),
\label{eq:device}
\end{equation}
its team update
\begin{equation}
w_i^{t,k+1} = (1 - \eta\lambda - \eta\gamma)\,w_i^{t,k} + \eta\gamma x^t + \lambda\eta\,\bar{\theta}_i^{t,k},
\label{eq:team}
\end{equation}
and its global update
\begin{equation}
x^{t+1} = (1 - \beta\gamma)x^t + \beta\gamma\bar{w}^t.
\label{eq:global}
\end{equation}
The parameter $\lambda$ appears in both \eqref{eq:device} and \eqref{eq:team}.
Section~\ref{sec:lambda-analysis} returns to this.

\section{Federated learning for intrusion detection}
\label{sec:fl-ids}

CFMD-i \parencite{cfmdi-2026} applies clustered federated learning to intrusion
detection across five datasets including CICIDS2017. It clusters by model
discrepancy, switching between a parameter-space measure and an output-space
Kullback-Leibler divergence according to a threshold that adapts to the current
spread of pairwise similarities. Clustering is triggered adaptively rather than
at a fixed round, on two conditions: the maximum inter-client model difference
must exceed a threshold, indicating enough heterogeneity to distinguish groups,
while the mean difference must stay below another, indicating shared structure
remains. Its contribution is communication efficiency, reporting a reduction of
over ninety-five per cent through parameter-difference transmission, adaptive
leapfrog communication and quantised error feedback.

Its client construction is relevant here. CFMD-i partitions traffic into attack
domains and assigns each participant one domain in its first scenario and
several in its second, so group structure exists by construction rather than
being imposed by a random draw.

ClusterFed \parencite{clusterfed-2025} uses self-supervised clustering with
balanced assignment via Sinkhorn-Knopp, which prevents cluster collapse. It
partitions clients by Dirichlet sampling at $\alpha = 0.1$.

Stones From Other Hills \parencite{soh-2025} studies self-labelled personalised
federated learning for IoT intrusion detection. Two aspects matter for this
project. It sweeps the degree of label skew, reporting FedAvg, a local baseline
and its own method at each level, and identifies an intermediate setting as the
usable balance: too much class overlap leaves nothing for personalisation to
exploit, too little prevents convergence. It also evaluates clients on classes
absent from their training data, so that adaptation to unseen traffic can be
measured.

Fed-ANIDS \parencite{fedanids-2023} takes an anomaly detection approach,
training autoencoders on normal traffic only and treating reconstruction error
as an intrusion score.

\section{Data quality in CICIDS2017}
\label{sec:data-quality}

CICIDS2017 \parencite{sharafaldin-2018} is the most used benchmark in this
area, and the literature documents several defects that inflate reported
results when unaddressed.

Flow rate columns contain infinities where a rate was computed over a
zero-length flow, and the dataset contains a substantial number of duplicate
rows. Reported practice is to drop both, or to impute the infinities with a
column median fitted inside the training fold. Identifier columns, addresses,
ports, timestamps and flow identifiers, must be removed before the feature
matrix is built, because a personalised model can otherwise score well by
recognising which client it is serving.

Percentile clipping is applied in several published pipelines to bound the
extreme values typical of flow statistics.

Class imbalance is severe. Benign traffic is roughly eighty-two per cent of the
capture, so accuracy is uninformative: a classifier that never raises an alert
scores above 0.81. Macro-averaged F1 weights each class equally and does not
use the true negative cell, which is why it separates a working detector from a
trivial one where accuracy does not.

\section{Gap}
\label{sec:gap}

Three gaps motivate the project.

PerMFL has not been evaluated on intrusion detection data. Its premise concerns
known team structures, and intrusion detection supplies them naturally, but no
published work tests it there.

The clustered methods that do address intrusion detection derive their
groupings and offer no convergence analysis; PerMFL offers the analysis but
assumes the grouping. Neither literature examines what the assumed grouping
actually contributes.

Where clustering methods report null or weak results, the partition is rarely
examined as a possible cause. Dirichlet partitioning produces a continuum of
client distributions rather than discrete groups, so a clustering algorithm has
nothing recoverable to find, and a null result measures the partition rather
than the method.
